\section{Objective and Methodology}

Building on the challenges and research gaps identified in the state of the art, this chapter defines the concrete goal of the thesis and the procedure chosen to achieve it. First, the overall objective and the core system functionalities are specified in the context of AR-assisted Programming by Demonstration (PbD). Then, the methodological approach is outlined, describing how the work progresses from requirements analysis to system design, implementation, and evaluation.

\subsection{Objective of the Thesis}

The objective of this thesis is to develop a proof-of-concept Augmented Reality (AR) application on the Apple Vision Pro using visionOS for low-latency six-degree-of-freedom (6D) pose estimation, serving as a basis for intuitive AR-based PbD systems. In contrast to classical teach-in and purely screen-based interfaces, the proposed system aims to tightly couple spatial perception, natural interaction, and pose feedback in a single mixed reality workflow.

The central idea is to close the loop between human demonstration and robot-perceivable scene understanding. To this end, the following core functionalities are defined:

\begin{itemize}
    \item \textbf{Object or region selection via gaze or hand gestures:} 
    Users should be able to select task-relevant objects or regions of interest (ROIs) directly in the spatial scene using natural input (gaze focus and pinch gestures). This replaces indirect selection via monitors or mouse input and is essential for PbD scenarios in which operators work in close proximity to the robot and the workpiece.

    \item \textbf{Low-latency 6D pose estimation via an external AI-based backend:}
    A dedicated backend service computes object poses from RGB(D) data and additional metadata. Offloading this step to an AI-based backend enables the use of state-of-the-art 6D pose estimators without exceeding the resource constraints of the AR device, while still targeting interactive update rates suitable for operator guidance.

    \item \textbf{AR visualization of estimated poses within the physical scene:}
    Estimated poses are rendered back into the live environment as virtual overlays (e.g., CAD models, axes, or wireframes). This provides direct visual feedback on pose quality, supports verification by the user, and forms the basis for later transfer of poses to robot controllers in PbD workflows.
\end{itemize}

Together, these functionalities are intended to demonstrate that a modern mixed reality headset can act as an intuitive front-end to an AI-based perception backend, creating a bridge between human demonstrations and machine-readable 6D object poses.

\subsection{Methodological Approach}

The methodological approach describes how the defined objective is pursued in a structured and traceable way. The work is organized into four phases that align with the subsequent chapters of the thesis.

% \textbf{Literature review and requirements definition} – Prior work on PbD, AR-based interaction, and 6D pose estimation is analyzed to identify relevant concepts, limitations, and design patterns. From this analysis, functional requirements (e.g., interaction modes, expected pose feedback, usability) and technical requirements (e.g., latency bounds, frame rates, data formats, and device constraints) for the prototype are derived.

\textbf{System architecture design} – Based on the requirements, the core components of the system are defined, including the AR front-end on visionOS, the backend pose estimation service, and the communication layer between them. Data flows and protocols are specified, with particular emphasis on how sensor input from the Vision Pro (RGB, depth, gaze, and gestures) is transformed into backend requests and how pose results are integrated into the AR scene.

\textbf{Application implementation} – The AR application is implemented on visionOS using ARKit, RealityKit, and SwiftUI. This phase covers the realization of gaze- and gesture-based ROI selection, the setup of network communication with the backend, and the rendering of 6D poses as spatial overlays. Performance-oriented design choices are applied to keep end-to-end latency low enough for interactive use.

\textbf{Evaluation} – Finally, the prototype is evaluated with respect to technical performance and user interaction. Technical metrics include latency, accuracy and robustness under different conditions (e.g., lighting, motion, partial occlusion), while usability is assessed through structured user tasks and subjective feedback. The results are used to discuss strengths, limitations, and opportunities for future extensions towards full robot programming by demonstration.
